\subsection{Round Brackets}\label{sec:Parenthesis}
Round Brackets (or Parentheses)\index{round brackets}\index{parentheses|see{round brackets}} are of course no whitespace characters – in the opposite, they are important syntax elements. But they trigger the use of whitespace, in particular blank spaces, and therefore, they are mentioned in this chapter.

We distinguish between \textit{arithmetical} round brackets and \textit{parameter} round brackets; the first is used in (mathematical or logical) expressions to order the terms, the latter is used with method or constructor calls or in lambdas.

The \bdq{Code Conventions for the \Java Programming Language}\autocite{SUN_CODE_CONVENTIONS:BlankSpaces}\index{Code Conventions for the Java Programming Language} does not make this distinction, and it handles them differently than described here (basically in the same way as described here as arithmetical round brackets).

\paragraph{Arithmetical Round Brackets}\index{arithmetical round brackets}
An opening arithmetical round bracket is always \textit{preceded} by a blank or another opening round bracket and \textit{never} followed by a blank. A closing arithmetical round bracket is \textit{never preceded} by a blank, but \textit{always followed} either by another closing round bracket, a blank, or a newline (or a semicolon, in case the statement ends).

This is a sample for arithmetical round brackets: \lstinline|((5 + 7) * 9 + 3) / 4|.

\paragraph{Parameter Round Brackets}\label{sec:ParameterParenthesis}\index{parameter round brackets}
An opening parameter round bracket is \textit{never preceded} by a blank and \textit{always followed} by one, while a closing parameter parenthesis is \textit{always preceded} by a blank. The only exception is for empty arguments lists where the opening parenthesis is immediately followed by the closing parenthesis.

Of course this is also valid for the declaration of methods and constructors.

Regarding to this rule, \lstinline|if|\index{if}, \lstinline|for|\index{for}, \lstinline|while|\index{while}, \lstinline|try|\index{try-with-resources} (with resources) and \lstinline|switch|\index{switch} are treated like they are functions.\footnote{And also the keywords \lstinline|catch|\index{try!catch} and \lstinline|synchronized|\index{synchronized}, but they do not allow complex terms inside their 'argument list'.}

Applying these rules can improve the readability of complex terms drastically, compared with the format suggested by the \ccjpl, as the examples below will illustrate.

\clearpage

\begin{tabular}{p{0.45\textwidth} p{0.45\textwidth}}
\hline
\lstinline[keepspaces=true]|final double d = sin( (a + b) / c );|
& 
\lstinline[keepspaces=true]|final double d = sin ((a + b) / c);|
\\
\hline
\lstinline[keepspaces=true]|System.out.println( "Finished!"  );|
&
\lstinline[keepspaces=true]|System.out.println ("Finished!");|
\\
\hline
\lstinline[keepspaces=true]|set( get() );|
&
\lstinline[keepspaces=true]|set (get ());|
\\
\hline
\lstinline[keepspaces=true]|System.out.println( sin( ((a + b) * (c + 4) + c)) / d ) );|
&
\lstinline[keepspaces=true]|System.out.println (sin (((a + b) * (c + 4) + c)) / d));|
\\
\hline
\lstinline[keepspaces=true]|if( a < b ) c = (d + e) / sin( f );|
&
\lstinline[keepspaces=true]|if (a < b) c = (d + e) / sin (f);|
\\
\hline
\lstinline[keepspaces=true]|while( ((line = reader.readLine()) != null) && (++i < capacity) ) processLine( line );|
&
\lstinline[keepspaces=true]|while (((line = reader.readLine ()) != null)~&& (++i < capacity)) processLine (line);|
\\
\hline
\end{tabular}

\keeplisting{3}
A special case are round brackets used with a cast:
\begin{lstlisting}
final Object o = "Text";
final var s = (String) o;
\end{lstlisting}

Here an opening round bracker is never followed by a blank, and a closing one is never preceded by a blank.

\subsubsection{Lambda Expressions}
According to the rules above, a lambda expression\index{lambda} with more than one argument should be written like this:
\keeplisting{3}
\begin{lstlisting}
final BiFunction f1 = ( a, b ) -> (a + b) * PI;
final BiFunction f2 = ( a, _ ) -> a * PI;
final BiFunction f3 = ( _, _ ) -> PI;
\end{lstlisting}

But it was found that it is easier to read without the blanks in the parameter list,  at least for single-letter parameter names or the unnamed variable\index{unnamed variable} (the underscore that is provided instead of the parameter):
\keeplisting{3}
\begin{lstlisting}
final BiFunction f1 = (a,b) -> (a + b) * PI;
final BiFunction f2 = (a,_) -> a * PI;
final BiFunction f3 = (_,_) -> PI;
\end{lstlisting}

So I suggest this style for lambda expressions.

%==============================================================================
% End of file!!
%==============================================================================
